\label{sec:results}

\subsection{Correlation Between Compactness and Driving Time}

Figure~\ref{fig:scatter} (described below; see replication materials for plot code)
shows the scatter of district-level mean constituent driving time against Polsby-Popper
compactness for all 432 included districts. The correlation is $r = -0.67$
(95\% CI: $[-0.72, -0.61]$, $p < 0.001$). The relationship is approximately log-linear:
a doubling of Polsby-Popper from 0.15 to 0.30 (a typical range in the enacted maps)
is associated with a reduction in mean driving time of approximately 11 minutes.

\textbf{Robustness to compactness metric}: The correlation is consistent across
alternative compactness metrics:

\begin{table}[ht]
\centering
\caption{Correlation Between Compactness and Mean Constituent Driving Time}
\label{tab:correlations}
\begin{tabular}{lrrr}
\toprule
Compactness Metric & All Districts ($n=432$) & Gerrymanders ($n=180$) & Neutral ($n=252$) \\
\midrule
Polsby-Popper & $-0.67$ & $-0.71$ & $-0.52$ \\
Reock & $-0.59$ & $-0.63$ & $-0.44$ \\
Schwartzberg & $-0.61$ & $-0.66$ & $-0.48$ \\
Neck Width Ratio & $-0.71$ & $-0.74$ & $-0.55$ \\
\bottomrule
\end{tabular}
\begin{tablenotes}
\small
\item All correlations are Pearson $r$, significant at $p < 0.001$.
  Neck Width Ratio is defined in Section~\ref{sec:mechanism}.
\end{tablenotes}
\end{table}

The neck width ratio (defined formally in Section~\ref{sec:mechanism}) shows a somewhat
stronger correlation than any of the standard compactness metrics, consistent with the
hypothesis that district necks are the geometric feature most directly responsible for
the compactness-distance relationship.

\subsection{State-Level Analysis}

At the state level, the pattern is consistent with the district-level finding but with
higher variance due to smaller sample sizes. Table~\ref{tab:states} reports the five
states with the largest and smallest enacted-vs-bisect driving time gaps.

\begin{table}[ht]
\centering
\caption{Largest and Smallest Enacted vs.\ \textsc{bisect} Driving Time Gaps}
\label{tab:states}
\begin{tabular}{lrrrl}
\toprule
State & Enacted Mean (min) & Bisect Mean (min) & Gap (min) & EG Classification \\
\midrule
\multicolumn{5}{l}{\textit{Largest gaps (enacted $\gg$ \textsc{bisect})}} \\
Maryland & 61.4 & 31.2 & 30.2 & Gerrymander ($-0.13$) \\
North Carolina & 56.8 & 33.7 & 23.1 & Gerrymander ($+0.11$) \\
Ohio & 54.2 & 33.4 & 20.8 & Gerrymander ($+0.12$) \\
Georgia & 52.7 & 34.1 & 18.6 & Gerrymander ($+0.09$) \\
Texas & 58.1 & 39.8 & 18.3 & Gerrymander ($+0.10$) \\
\midrule
\multicolumn{5}{l}{\textit{Smallest gaps (enacted $\approx$ \textsc{bisect})}} \\
Iowa & 38.2 & 37.1 & 1.1 & Neutral ($+0.02$) \\
California & 35.4 & 34.8 & 0.6 & Commission ($-0.01$) \\
Colorado & 33.1 & 32.9 & 0.2 & Commission ($+0.02$) \\
Minnesota & 36.8 & 36.4 & 0.4 & Court-drawn ($-0.01$) \\
Virginia & 37.2 & 37.9 & $-0.7$ & Commission ($+0.03$) \\
\bottomrule
\end{tabular}
\begin{tablenotes}
\small
\item EG = Efficiency Gap from enacted 2022 map. Gerrymander threshold: $|\text{EG}| > 0.08$.
\end{tablenotes}
\end{table}

Three observations stand out:

\begin{enumerate}
\item The five states with the largest enacted-vs-bisect gaps are all classified as
  partisan gerrymanders, with the EG sign indicating which party benefits. Maryland
  ($\text{EG} = -0.13$) benefits Democrats; the remaining four benefit Republicans.
  This is consistent with the prediction that gerrymandering, not compact algorithmic
  planning, drives the elevated driving times.

\item Commission and court-drawn states show near-zero gaps, and in one case (Virginia)
  the enacted plan is marginally better than the \textsc{bisect} plan---within noise.
  This provides a placebo test: in states without gerrymandering, compact criteria are
  already largely satisfied by the enacting body.

\item Texas shows a large gap ($+18.3$ minutes) despite having large average district
  area (Texas is geographically large). This suggests that in Texas, the driving time
  penalty from gerrymandering operates not primarily through district area but through
  the geometric ``necks'' that connect distant communities.
\end{enumerate}

\subsection{Regression Analysis}

We estimate the following OLS regression:

\begin{equation}
\bar{t}_d = \alpha + \beta \cdot \text{PP}_d + \mathbf{X}_d'\gamma + \varepsilon_d
\label{eq:regression}
\end{equation}

where $\bar{t}_d$ is the mean constituent driving time in minutes for district $d$,
$\text{PP}_d$ is the Polsby-Popper score, and $\mathbf{X}_d$ is a vector of controls
including log district area (km$^2$), log district population (constant at 760,000 under
one-person-one-vote), state fixed effects, and urban fraction (fraction of district
population classified as urban by the Census Bureau).

\begin{table}[ht]
\centering
\caption{OLS Regression: Polsby-Popper and Mean Constituent Driving Time}
\label{tab:regression}
\begin{tabular}{lrrrr}
\toprule
 & (1) & (2) & (3) & (4) \\
 & Bivariate & $+$Area & $+$Urban & $+$State FE \\
\midrule
Polsby-Popper & $-58.3^{***}$ & $-41.2^{***}$ & $-33.7^{***}$ & $-22.4^{***}$ \\
 & $(4.1)$ & $(5.2)$ & $(5.8)$ & $(6.7)$ \\
Log District Area & & $8.9^{***}$ & $6.4^{***}$ & $4.1^{**}$ \\
 & & $(1.1)$ & $(1.3)$ & $(1.8)$ \\
Urban Fraction & & & $-18.6^{***}$ & $-14.3^{***}$ \\
 & & & $(2.9)$ & $(3.1)$ \\
State Fixed Effects & No & No & No & Yes \\
\midrule
$R^2$ & 0.449 & 0.521 & 0.607 & 0.721 \\
$N$ & 432 & 432 & 432 & 432 \\
\bottomrule
\end{tabular}
\begin{tablenotes}
\small
\item $^{***} p < 0.01$, $^{**} p < 0.05$. Standard errors clustered by state.
  Dependent variable: mean constituent driving time (minutes).
\end{tablenotes}
\end{table}

The Polsby-Popper coefficient is negative and significant across all specifications.
In the most demanding specification with state fixed effects (column 4), a one standard
deviation increase in Polsby-Popper ($\sigma = 0.11$) is associated with a reduction
in mean constituent driving time of $0.11 \times 22.4 = 2.5$ minutes within state.
The within-state effect is attenuated relative to the cross-state effect (column 1:
$0.11 \times 58.3 = 6.4$ minutes) because state fixed effects absorb variation in
geography and state-specific redistricting practices that are correlated with
both compactness and office location choices.
